\chapter{Constexpr deep dive: how the DSL computes at build time}
\label{ch:constexpr-deep}

\begin{aside}
\code{constexpr} and \code{consteval} are the engines of the
DSL. Most users see them as decorations; behind the scenes they
are full programming languages running inside the compiler. This
chapter takes a deep look at what is computable at compile time
in C++26, where the limits are, and how the DSL uses every
available feature.
\end{aside}

\section{What runs at compile time}

A \code{constexpr} function may run at compile time when its
arguments are constant expressions. A \code{consteval} function
\emph{must} run at compile time --- the compiler errors out if a
runtime invocation is attempted.

The set of operations allowed inside \code{constexpr} has
grown dramatically:

\begin{itemize}[leftmargin=*]
  \item C++11: arithmetic, simple branches, single return.
  \item C++14: loops, mutation of locals, multiple returns.
  \item C++17: lambda invocation, \code{if constexpr}.
  \item C++20: dynamic allocation (transient), virtual functions,
    \code{try/catch} (limited).
  \item C++23: \code{static\_assert} of non-literal types,
    \code{constexpr} std::optional, std::variant.
  \item C++26: even more relaxations; reflection in the pipeline.
\end{itemize}

The DSL uses C++20 and later. The compiler can run real
programs --- not just constant folding --- to compute its NTTPs.

\section{A worked example: parsing a hex colour at compile time}

\code{Fg<"\#FF6432">} requires the compiler to parse the hex
literal and convert it to three integers. The code:

\begin{lstlisting}
constexpr uint8_t hex_digit(char c) {
    if (c >= '0' && c <= '9') return c - '0';
    if (c >= 'a' && c <= 'f') return c - 'a' + 10;
    if (c >= 'A' && c <= 'F') return c - 'A' + 10;
    throw std::invalid_argument(\"not a hex digit\");
}

constexpr std::array<uint8_t, 3> parse_hex(std::string_view s) {
    if (s.size() != 7 || s[0] != '#')
        throw std::invalid_argument(\"expected #RRGGBB\");
    return {
        static_cast<uint8_t>(hex_digit(s[1]) * 16 + hex_digit(s[2])),
        static_cast<uint8_t>(hex_digit(s[3]) * 16 + hex_digit(s[4])),
        static_cast<uint8_t>(hex_digit(s[5]) * 16 + hex_digit(s[6]))
    };
}

template <Str s>
constexpr auto fg = []() {
    constexpr auto rgb = parse_hex(s.view());
    return Fg<rgb[0], rgb[1], rgb[2]>;
}();
\end{lstlisting}

Several things happen here:

\begin{itemize}[leftmargin=*]
  \item \code{parse\_hex} is \code{constexpr}. It throws on
    malformed input.
  \item The thrown exception, in a \code{constexpr} context, is
    a compile-time error. The compiler reports it as
    ``call to non-constexpr function'' or ``thrown exception''.
  \item \code{template <Str s>} takes the literal as an NTTP;
    the immediately-invoked lambda runs at compile time and
    instantiates \code{Fg} with the parsed values.
\end{itemize}

\subsection*{What the compiler sees}

When you write \code{fg<"\#FF6432">}, the compiler:

\begin{enumerate}[leftmargin=*]
  \item Builds a \code{Str<8>} from the literal.
  \item Substitutes into the template.
  \item Evaluates the lambda body: calls \code{parse\_hex}.
  \item \code{parse\_hex} returns \code{[255, 100, 50]}.
  \item Instantiates \code{Fg<255, 100, 50>}.
  \item Caches the result; \code{fg<"\#FF6432">} is now this
    instance.
\end{enumerate}

All of that happens before any runtime code runs.

\section{Throwing in constexpr}

A throw inside \code{constexpr} is interesting:

\begin{itemize}[leftmargin=*]
  \item If the function is called at runtime, throwing is
    normal: the exception propagates.
  \item If the function is called at compile time, the throw
    is a compile error: ``a thrown exception is not a constant
    expression''.
\end{itemize}

This is the mechanism by which \code{parse\_hex} validates input.
You get runtime exceptions in runtime context and compile errors
in constexpr context, from the same source.

\begin{idea}
This dual nature is what makes \code{constexpr} so powerful for
DSLs. A single function can validate its inputs at \emph{either}
compile time or runtime, depending on how it is called. The DSL
is just a function called at compile time with literal
arguments.
\end{idea}

\section{consteval: compile-time only}

When you want a function that must never run at runtime, use
\code{consteval}:

\begin{lstlisting}
consteval CTStyle merge_styles(CTStyle a, CTStyle b) {
    return CTStyle{
        .bold = a.bold || b.bold,
        .italic = a.italic || b.italic,
        // ...
    };
}
\end{lstlisting}

The compiler errors if you try to call \code{merge\_styles}
with a runtime argument. This is what gates the DSL's pipe
operators: they call \code{merge\_styles}, which insists on
constexpr inputs, which forces the pipe to be a compile-time
expression.

\section{constinit: ensuring static-init order}

\begin{lstlisting}
constinit Theme default_theme = make_dark_theme();
\end{lstlisting}

\code{constinit} guarantees the variable is initialised at
compile time, not via dynamic initialisation. This avoids the
static-init order fiasco: \code{constinit} variables are ready
before \code{main} runs.

\maya{} uses \code{constinit} for theme defaults, palette
tables, and other startup-critical values.

\section{constexpr std::vector and std::string}

C++20 made \code{std::vector} and \code{std::string} usable in
\code{constexpr}. The compiler implements a tiny heap inside the
constant evaluator:

\begin{lstlisting}
consteval std::vector<int> build_table() {
    std::vector<int> v;
    for (int i = 0; i < 100; ++i) v.push_back(i * i);
    return v;
}

constexpr auto squares = build_table();
\end{lstlisting}

The catch: the vector must not escape the constant evaluator.
You cannot return a \code{constexpr std::vector} to runtime;
the heap allocation is freed when constexpr evaluation ends.

The workaround: copy into a \code{std::array}:

\begin{lstlisting}
consteval auto fixed_squares() {
    auto v = build_table();
    std::array<int, 100> a{};
    for (size_t i = 0; i < 100; ++i) a[i] = v[i];
    return a;
}

constexpr auto squares = fixed_squares();
\end{lstlisting}

\maya{} uses this pattern for the Unicode width table: build at
compile time, store as \code{constexpr std::array}.

\section{if constexpr: compile-time branches}

\begin{lstlisting}
template <class T>
auto encode(T value) {
    if constexpr (std::is_integral_v<T>) {
        return std::to_string(value);
    } else if constexpr (std::is_floating_point_v<T>) {
        return std::format(\"{:f}\", value);
    } else {
        return std::string{value};
    }
}
\end{lstlisting}

Branches whose condition is \code{constexpr} are evaluated at
compile time; only the chosen branch is instantiated. The other
branches do not have to be well-formed for the unselected type.

\maya{} uses \code{if constexpr} for platform-specific code,
SIMD level selection, and template specialisation without
explicit overloading.

\section{requires expressions in constexpr}

Inside a function, you can use \code{requires} to check
properties:

\begin{lstlisting}
template <class T>
auto format(T value) {
    if constexpr (requires { value.to_string(); }) {
        return value.to_string();
    } else if constexpr (requires { std::to_string(value); }) {
        return std::to_string(value);
    } else {
        return std::string{\"(no format)\"};
    }
}
\end{lstlisting}

The requires-expression is a compile-time predicate. The
branches are selected based on what types support what
operations.

\section{The limits of compile-time computation}

What you cannot do at compile time:

\begin{itemize}[leftmargin=*]
  \item I/O. No file reads, no network, no stdout.
  \item Reading the system clock.
  \item Random numbers (other than constexpr-friendly PRNGs).
  \item Reflection on names (until C++26 reflection).
  \item Unbounded recursion (the compiler has a depth limit,
    typically 1024).
\end{itemize}

The DSL stays within these limits. Every literal it processes,
every style it merges, every box it configures --- all done
without I/O or randomness.

\section{Compile-time vs runtime: who wins?}

A computation that fits in either context has a choice.
Compile-time advantages:

\begin{itemize}[leftmargin=*]
  \item Zero runtime cost.
  \item Errors surface at build, not in production.
  \item Result is a known value; the optimiser sees it.
\end{itemize}

Runtime advantages:

\begin{itemize}[leftmargin=*]
  \item Faster builds: not every value needs to be a template
    parameter.
  \item More flexible: runtime input.
  \item Debuggable.
\end{itemize}

The DSL's principle: shape at compile time, content at runtime.
Box configurations, style flags, sizes --- compile time. Text
contents, list lengths, model values --- runtime.

\section{Build time as a cost}

A complex DSL expression compiles slowly. \maya{}'s typical
build:

\begin{itemize}[leftmargin=*]
  \item A small app (1000 lines): 2--5 seconds.
  \item A medium app (10\,000 lines): 30--60 seconds.
  \item A large app (100\,000 lines): minutes.
\end{itemize}

Most of the time is spent in template instantiation. If you
care about build times, use the runtime variants
(\code{text(\ldots)}, \code{box(\ldots)}) for the bulk of your
UI and reserve the compile-time DSL for hot or critical paths.

\section*{Workshop}

\begin{enumerate}[leftmargin=*]
  \item Write a \code{constexpr} function that computes the
    factorial. Call it at compile time and at runtime; observe
    both work.
  \item Add a \code{static\_assert} that fails. Read the
    compiler error. Now make it pass.
  \item Read \file{maya/include/maya/dsl.hpp} and find every
    \code{consteval} function. Each is a piece of the DSL's
    compile-time computation.
\end{enumerate}

\begin{recap}
\code{constexpr} and \code{consteval} are the engines of
compile-time computation. The DSL uses them to parse literals,
merge styles, validate configurations, and build static lookup
tables. The compiler runs real programs to compute the types
your runtime sees. The cost is build time; the payoff is zero
runtime overhead and compile-time error detection.
\end{recap}
